\section{Spherical experiments with a spectral gap}\label{sec:nonabelian56}
The next result concerns a full vector orbit, not an abelian factor of a nonabelian system. Write $\sigma$ for normalized area on $S^{d-1}$, where $d\ge2$. Let a finite orthogonal alphabet contain the identity, and let $\mu$ be a probability law on its letters. Assume that its averaging operator
\[
                  Pf(x)=\E_{R\sim\mu}f(Rx)
\]
satisfies
\begin{equation}\label{eq:spectralgap56}
        \|P^n f\|_{L^2(\sigma)}\le\lambda^n\|f\|_{L^2(\sigma)}
        \quad\left(\int f\,d\sigma=0\right),\qquad 0<\lambda<1.
\end{equation}
The physical target is $F(g)=\rho g x_0$ for a fixed $x_0\in S^{d-1}$, with output set the closed Euclidean unit ball and error measured in Euclidean norm.

\begin{theorem}[A spherical width law up to a logarithm]\label{thm:spherical56}
Assume \eqref{eq:spectralgap56}, $0<\rho<1$, and $0\le\varepsilon<\rho$. There are constants $c,C>0$ such that
\begin{equation}\label{eq:sphericalrate56}
 c\left(\frac{N}{\log(N+2)}\right)^{(d-1)/2}
       \le W_{N,\varepsilon}\le C(N+1)^{(d-1)/2}.
\end{equation}
The upper realization is exact. Every error-$\varepsilon$ realization has at most
\[
                      C k^{2/(d-1)}\log(k+2)
\]
positive cuts of width at most $k$, with other widths unrestricted. The constants may depend on $d,\lambda,\rho,\varepsilon$, but not on the horizon or the realization. At unit amplitude the same bounds hold for every fixed $0<\varepsilon<1$; exact unit amplitude is not included in the upper assertion.
\end{theorem}

\begin{lemma}[Spherical quantization and executable contraction]\label{lem:sphericalgap56}
Under \eqref{eq:spectralgap56}, there are $c_d>0$ and integers $B_k\le C_{d,\lambda}\log(k+2)$ such that
\[
 \E_{w\sim\mu^{B_k}}\max_{i\le k}\langle v_i,u_w u\rangle
                  \le1-c_d k^{-2/(d-1)}
\]
for all unit vectors $u,v_1,\ldots,v_k$. The defect may be decreased so that it lies in $(0,1)$. When $d=3$ one may take defect $1/(2k)$ and
\begin{equation}\label{eq:sphere3constants56}
                  B_k=\left\lceil\frac{\log(64k^2)}{-\log\lambda}\right\rceil.
\end{equation}
\end{lemma}
\begin{proof}
Normalized spherical cap measure obeys
\[
       c_d h^{d-1}\le\sigma\{x:\|x-u\|_2<h\}\le C_d h^{d-1}
                     \qquad(0<h\le1),
\]
uniformly in $u$. This follows by writing area in polar angle as a constant times $\sin^{d-2}\theta\,d\theta$ and using $h=2\sin(\theta/2)$; for $d=2$ the same formula is the arc-length formula. For $f(x)=\max_i\langle v_i,x\rangle$, choose $h=c k^{-1/(d-1)}\le1$ so that the union of the $k$ radius-$h$ caps has area at most $1/2$. Outside them, $1-f(x)\ge h^2/2$. Thus
\begin{equation}\label{eq:haarsphere56}
                         \int f\,d\sigma\le1-a_d k^{-2/(d-1)}.
\end{equation}
The function $f$ is $1$-Lipschitz for chordal distance and bounded by one in absolute value.

Let $K_\delta$ average over the chordal cap of radius $\delta\le1$. It commutes with all rotations, fixes constants, and satisfies $\|K_\delta f-f\|_\infty\le\delta$. The cap density at a center has $L^2$ norm at most $C_d\delta^{-(d-1)/2}$. Hence \eqref{eq:spectralgap56} and Cauchy--Schwarz imply, uniformly in the center,
\[
 \left|P^B f(u)-\int f\,d\sigma\right|
 \le\delta+C'_d\delta^{-(d-1)/2}\lambda^B.
\]
Take $\delta$ to be a fixed sufficiently small multiple of $k^{-2/(d-1)}$ and then $B=O_{d,\lambda}(\log(k+2))$ so that the right side is at most half the defect in \eqref{eq:haarsphere56}. The law of the $B$ successive letters is a finite executable word law, proving the general assertion.

For $d=3$, a chordal cap of radius $h\le2$ has normalized area $h^2/4$. The union bound gives $\Pr(1-f<t)\le kt/2$, so integration for $0\le t\le2/k$ gives $\int(1-f)\,d\sigma\ge1/k$. A radius-$\delta$ cap density has $L^2$ norm $2/\delta$, and $\|f-\int f\|_2\le2$. Thus the mixing error is at most $\delta+4\lambda^B/\delta$. Choose $\delta=1/(4k)$ and \eqref{eq:sphere3constants56}; the error is at most $1/(2k)$.
\end{proof}

\begin{proof}[Proof of Theorem~\ref{thm:spherical56}]
Let $Z=u_wx_0$ and let $D(S)$ be the terminal decoder. For every deterministic word, correctness and $\|Z\|_2=1$ give
\[
           \langle\E[D(S)\mid w],Z\rangle\ge\rho-\varepsilon.
\]
Under any external word law this implies
\begin{equation}\label{eq:sphericalresidual56}
       \rho-\varepsilon\le\E\langle D(S),Z\rangle
             \le\E\|\E[Z\mid S]\|_2,
\end{equation}
since the decoder lies in the unit ball. Lemma~\ref{lem:centroid54}, applied with the executable law of Lemma~\ref{lem:sphericalgap56}, contracts the last conditional norm by $1-a_k$ at each selected narrow endpoint, where $a_k=c_d k^{-2/(d-1)}\in(0,1)$. Initially the norm is one. Fixed identity-letter gaps cannot increase it, even though their stochastic rows may process the register. After $J$ selected endpoints, \eqref{eq:sphericalresidual56} gives
\[
            J[-\log(1-a_k)]\le\log\frac1{\rho-\varepsilon}.
\]
Greedy spacing by $B_k$ as in the occupation proof discards at most $B_k-1$ initial cuts and charges at most $B_k$ eligible cuts to each selected endpoint. Therefore the occupation count is at most
\begin{equation}\label{eq:sphericaloccupation56}
 B_k\left(1+a_k^{-1}\log\frac1{\rho-\varepsilon}\right).
\end{equation}
If every register has width at most $k$, this bounds $N$. Inverting the bound gives the lower order in \eqref{eq:sphericalrate56}. Explicitly, when $k\le(N+1)^{d-1}$ its logarithm is $O_d(\log(N+2))$, and otherwise the asserted lower bound is automatic.

For the upper bound, choose a maximal $\delta$-separated set of vertices on $S^{d-1}$ containing $x_0$, with $0<\delta\le1/2$. The disjoint caps of radius $\delta/2$ show that it has at most $C_d\delta^{-(d-1)}$ vertices. Maximality makes it a $\delta$-net. Its convex hull $Q$ satisfies
\begin{equation}\label{eq:sphericalpolytope56}
                r\mathbb B^d\subset Q\subset\mathbb B^d,
                         \qquad r=1-\delta^2/2,
\end{equation}
because for every unit direction $u$ some vertex $v$ has $\langle u,v\rangle\ge1-\delta^2/2$. The inclusion follows from the supporting half-space characterization of a closed convex body. For each command $R$ and vertex $v_i$, the point $rRv_i$ belongs to $Q$. Choose barycentric coordinates of that point as the $i$th row of one common stochastic command matrix. The identity command is treated in the same way, with the same factor $r$.

Initialize at the vertex $x_0$ and decode vertex $v_i$ as $(\rho/r^N)v_i$. Induction on the row means gives output exactly $\rho u_wx_0$ on every length-$N$ word. The decoder is legal if $r^N\ge\rho$. Choose
\[
                 \delta^2=\min\left(\tfrac14,
                                      \frac{-\log\rho}{N+1}\right).
\]
Since $\log(1-\delta^2/2)\ge-\delta^2$, the required inequality holds. The number of vertices is $O_{d,\rho}((N+1)^{(d-1)/2})$. Both rows and decoder may depend on the horizon, as permitted by the clocked model.

For $\rho=1$ and $\varepsilon>0$, the same lower proof applies. An exact upper realization at amplitude $\rho'=1-\varepsilon/2$ differs from the unit-amplitude target by $\varepsilon/2$ in Euclidean norm, proving the stated upper bound. There is no assertion of exactness at unit amplitude.
\end{proof}

\begin{corollary}[A rational noncommuting alphabet on $\mathrm{SO}(3)$]\label{cor:SO356}
Let the alphabet consist of the identity, the matrices
\[
 R_x=\begin{pmatrix}1&0&0\\0&3/5&-4/5\\0&4/5&3/5\end{pmatrix},
 \qquad
 R_z=\begin{pmatrix}3/5&-4/5&0\\4/5&3/5&0\\0&0&1\end{pmatrix},
\]
and their inverses. For $F(g)=\rho g e_3$, with $0<\rho<1$ and $0\le\varepsilon<\rho$, the bounds are
\[
             c\frac{N}{\log(N+2)}\le W_{N,\varepsilon}\le C(N+1).
\]
At error $\varepsilon\ge\rho$ a single zero decoder suffices. Every subcritical realization has at most $Ck\log(k+2)$ positive cuts of width at most $k$.
\end{corollary}
\begin{proof}
The angle $\theta$ of both rotations has $2\cos\theta=6/5$. It is not a rational multiple of $2\pi$: otherwise $2\cos\theta$ would be a rational algebraic integer, hence an integer. The closures of the two cyclic groups are the full circle subgroups about their distinct axes. Their Lie algebras and bracket span $\mathfrak{so}(3)$, so the generated compact closure is $\mathrm{SO}(3)$. The matrices do not commute.

Under the double cover $\mathrm{SU}(2)\to\mathrm{SO}(3)$, their unit-quaternion lifts are $(2+i)/\sqrt5$ and $(2+k)/\sqrt5$. They have algebraic matrix entries and generate a dense subgroup of $\mathrm{SU}(2)$: its closed image is all of $\mathrm{SO}(3)$, so its Lie algebra has dimension three. The spectral-gap theorem of Bourgain--Gamburd \cite{BG2012} applies. Adding the identity with positive mass makes the symmetric averaging operator have norm strictly less than one on mean-zero $L^2$. The spherical representation is a subrepresentation of the regular representation, so the uniform law on the five displayed letters satisfies \eqref{eq:spectralgap56} on $S^2$. Apply Theorem~\ref{thm:spherical56}. The radius of the full sphere $\rho S^2$ inside the unit ball is $\rho$, attained at zero.
\end{proof}
The spectral gap in this corollary is an imported theorem, not a numerical assertion extracted from finite orbit checks. This retained block estimate has a logarithmic loss. Theorem~\ref{thm:sharpoccupation61} supplies the matching linear lower order under the same spectral-gap hypothesis. It does not compute a numerical gap for this particular alphabet; Theorem~\ref{thm:effectivelps61} gives a separate alphabet with explicit constants.
